% !PRINTED 342
\weberpage{342}
Sind $\alpha, \beta, \gamma \ldots$ beliebige reelle oder complexe Grössen, so jedoch, dass der absolute Werth von $\alpha$ \gesperrt{grösser} ist, als der absolute Werth aller übrigen, dass also die absoluten Werthe der Brüche $\beta:\alpha, \gamma:\alpha \ldots$ echte Brüche sind, so ist
\weq{1}{
\frac{\alpha^m + \beta^m + \gamma^m + \cdots}{\alpha^{m-1}+\beta^{m-1}+\gamma^{m-1}+\cdots}
= \alpha\,\frac{1+\left(\dfrac{\beta}{\alpha}\right)^m+\left(\dfrac{\gamma}{\alpha}\right)^m+\cdots}{1+\left(\dfrac{\beta}{\alpha}\right)^{m-1}+\left(\dfrac{\gamma}{\alpha}\right)^{m-1}+\cdots},
}
und je grösser $m$ wird, um so mehr wird sich dieser Ausdruck, wie die zweite Darstellung zeigt, der Grenze $\alpha$ nähern.

Sind $\alpha, \beta, \gamma \ldots$ die Wurzeln einer algebraischen Gleichung, so ist die linke Seite von (1) der Quotient der $m^{\mathrm{ten}}$ und $(m-1)^{\mathrm{ten}}$ Potenzsumme, und wir erhalten also den Satz:

\begin{center}
\gesperrt{Der Quotient der $m^{\mathrm{ten}}$ und $(m-1)^{\mathrm{ten}}$ Potenzsumme nähert sich mit wachsendem $m$ der absolut grössten unter den Wurzeln.}
\end{center}

Nimmt man $m$ negativ an, und setzt $\alpha$ absolut \gesperrt{kleiner} als $\beta, \gamma \ldots$ voraus, so folgt auf die gleiche Weise:

\begin{center}
\gesperrt{Der Quotient der $-m^{\mathrm{ten}}$ und $-(m+1)^{\mathrm{ten}}$ Potenzsumme nähert sich mit wachsendem $m$ der absolut kleinsten unter den Wurzeln.}
\end{center}

Da man die Potenzsummen als symmetrische Functionen durch die Coëfficienten berechnen kann, so braucht man nur ein hinlänglich grosses $m$ zu nehmen, um einen angenäherten Werth der absolut grössten und absolut kleinsten Wurzel zu erhalten.

Wenn aber unter den Wurzeln der Gleichung solche vorkommen, die denselben absoluten Werth haben, und dieser der grösste oder der kleinste ist, so ist diese Methode nicht anwendbar, also z.\,B. nicht auf reelle Gleichungen, bei denen ein Paar imaginärer Wurzeln vorkommt, das die reellen an absolutem Werth übertrifft.

Wenn aber die grösste Wurzel zwar einzeln vorkommt, aber die nächstfolgende nicht viel übertrifft, so wird das Verfahren nur langsam einige Genauigkeit geben.

\gesperrt{Jacobi} hat die \gesperrt{Bernoulli}'sche Methode nach einer Richtung ergänzt\footnote{Observatiunculae etc. Werke, Bd. 3, S. 280.}.
